\section{Introduction}\label{sec:intro}

In October 2018, Google released BERT, a bidirectional Transformer encoder that
set the state of the art on eleven natural-language understanding tasks and
pushed the GLUE benchmark to 80.5\%~\cite{devlin2019bert}. BERT could not write
a coherent paragraph; it was a \emph{reading} model, fine-tuned per task. In
July 2026, Anthropic released Claude Opus~5 and reported 42/42 under its
post-contest, no-tools grading protocol for the 2026 International Mathematical
Olympiad problems; ARC Prize independently measured a roughly fourfold jump on
ARC-AGI-3~\cite{anthropic2026opus5card,arcprize2026opus5}. Between those
two releases sit eight years in which the field's dominant paradigm changed at
least four times: from task-specific fine-tuning, to few-shot prompting of
large pretrained decoders, to instruction tuning and human-feedback alignment,
to reasoning models and tool-using agents.

We ask three related questions. How did the field move from BERT to the
2026 frontier of GPT-5.6, Claude Fable~5 and Opus~5, Kimi~K3, and their
contemporaries? How quickly did capability and cost efficiency improve when
measured from public benchmark and pricing series rather than impressions?
And how much more performance can a practitioner obtain from fixed model
weights by spending inference compute carefully, calibrating reasoning effort,
estimating confidence, and routing work across models?

\paragraph{Contributions.}
\begin{enumerate}
    \item A consolidated, sourced timeline of language-model capability from
    2018 to 2026 across three benchmark eras (Section~\ref{sec:history},
    Table~\ref{tab:timeline}).
    \item Trend fits showing ${\sim}5.8\times$ annual growth in the odds of
    solving SWE-bench Verified tasks, alongside the saturation of
    knowledge benchmarks such as MMLU (Section~\ref{sec:capability}).
    \item A cost analysis showing a ${\sim}60\times$ decline in input-token
    prices from GPT-3 (2020) to GPT-5.6 Luna (2026), and a demonstration that
    today's budget tier matches the flagship of one quarter ago on most
    agentic/professional benchmarks (Section~\ref{sec:cost}).
    \item Evidence that the 2026 frontier is task-fragmented, and that simple
    model routing recovers the per-task optimum
    (Section~\ref{sec:specialization}).
    \item A locked development/evaluation study of self-consistency for a
    small open model, followed by an exploratory confidence model that ranks
    predictions for triage. Complete traces, paired uncertainty, and
    reasoning-effort analyses accompany the results
    (Section~\ref{sec:improvements}).
\end{enumerate}

\section{Eight years in four paradigms}\label{sec:history}

\subsection{The encoder era: BERT and the fine-tuning paradigm (2018--2019)}
BERT~\cite{devlin2019bert} established the pretrain-then-fine-tune recipe:
masked-language-model pretraining on unlabeled text, followed by a small
task-specific head. RoBERTa~\cite{liu2019roberta} showed the recipe was
undertrained rather than wrong, reaching 88.5 on GLUE with more data and longer
training; T5~\cite{raffel2020t5} unified all tasks as text-to-text and reached
90.3 on SuperGLUE. GPT-2~\cite{radford2019gpt2} demonstrated that a
left-to-right decoder trained on enough web text acquired tasks
\emph{without} fine-tuning, but its zero-shot quality was not yet competitive.
The era's signature: capability lived in task-specific heads, and benchmarks
(GLUE, SuperGLUE, SQuAD) saturated within roughly two years of introduction.

\subsection{Scale and in-context learning: GPT-3 (2020--2021)}
GPT-3~\cite{brown2020gpt3} scaled the decoder recipe to 175B parameters and
showed that few-shot prompting --- a handful of examples in the context window
--- could substitute for gradient updates on many tasks. Its 43.9\%
five-shot MMLU~\cite{hendrycks2021mmlu} looks modest today, but the conceptual
shift was permanent: capability could be \emph{elicited} from a frozen general
model. The limits were equally clear: GPT-3 few-shot scored 71.8 on SuperGLUE,
well below fine-tuned systems, and its unaligned samples were unreliable.

\subsection{Alignment and instruction following (2022--2023)}
InstructGPT~\cite{ouyang2022instructgpt} applied reinforcement learning from
human feedback (RLHF) to make models follow instructions; ChatGPT (November
2022) packaged this for general users. GPT-4 (March 2023) jumped to 86.4\% on
MMLU~\cite{openai2023gpt4}, near the estimated human-expert level, and
established the modern pattern of a capability report accompanying release.
Meta's Llama series opened weights and created an open ecosystem; Claude and
Gemini established a multi-vendor frontier.

\subsection{Reasoning and agents (2024--2026)}
Test-time compute changed the scaling target. OpenAI's o1 (December 2024)
showed that longer, trained reasoning traces and additional inference tokens
could lift math and coding performance; its 91.8\% MMLU score effectively
saturated that benchmark~\cite{benchlm_mmlu}. Agency then changed the task
itself. SWE-bench Verified~\cite{jimenez2024swebench,openai2024swebenchverified}
made ``resolve a real GitHub issue'' the standard hard problem. Scores on
independent harnesses climbed from 49\% for Claude 3.5 Sonnet in October 2024
to 97\% for Claude Opus~5 in July 2026~\cite{vals2026swebench}. By mid-2026
the flagship
releases --- Claude Fable~5 (July 1), GPT-5.6 Sol/Terra/Luna (July 9), Kimi~K3
(July 16), Claude Opus~5 (July 24) --- were evaluated primarily on agentic,
long-horizon, and economically-weighted tasks rather than static question
answering~\cite{anthropic2026fable5,openai2026gpt56,moonshot2026k3,anthropic2026opus5}.
Kimi~K3 additionally marked the near-closure of the open--closed gap: a
2.8-trillion-parameter open-weight mixture-of-experts model competitive with
the best proprietary systems~\cite{moonshot2026k3}.

Table~\ref{tab:timeline} summarizes the arc, and Figure~\ref{fig:swe} shows
the agentic benchmark series at its current endpoint.

\begin{table}[t]
\centering
\caption{Selected milestones, 2018--2026. Scores are as originally reported by
the cited source; benchmark variants are noted.}\label{tab:timeline}
\small
\begin{tabular}{@{}llllr@{}}
\toprule
Date & Model & Org & Milestone metric & Score \\
\midrule
2018-10 & BERT-Large & Google & GLUE & 80.5 \\
2019-07 & RoBERTa & Meta & GLUE & 88.5 \\
2019-10 & T5-11B & Google & SuperGLUE & 90.3 \\
2020-05 & GPT-3 & OpenAI & MMLU (5-shot) & 43.9 \\
2023-03 & GPT-4 & OpenAI & MMLU (5-shot) & 86.4 \\
2024-05 & GPT-4o & OpenAI & MMLU & 88.7 \\
2024-12 & o1 & OpenAI & MMLU & 91.8 \\
2024-10 & Claude 3.5 Sonnet & Anthropic & SWE-bench Ver. & 49.0 \\
2025-05 & Claude Opus 4 & Anthropic & SWE-bench Ver. & 72.5 \\
2025-08 & GPT-5 & OpenAI & SWE-bench Ver. & 74.9 \\
2025-11 & Claude Opus 4.5 & Anthropic & SWE-bench Ver. & 76.8 \\
2026-04 & GPT-5.5 & OpenAI & SWE-bench Ver. (vendor) & 85.1 \\
2026-07 & Claude Fable 5 & Anthropic & SWE-bench Ver. (indep.) & 95.0 \\
2026-07 & GPT-5.6 Sol & OpenAI & SWE-bench Ver. (indep.) & 96.2 \\
2026-07 & Kimi K3 & Moonshot & SWE-bench Ver. (indep.) & 93.4 \\
2026-07 & Claude Opus 5 & Anthropic & SWE-bench Ver. (indep.) & 97.0 \\
\bottomrule
\end{tabular}
\end{table}
